\documentclass[11pt]{article}

\usepackage[margin=0.85in, top=0.75in, bottom=0.75in]{geometry}
\usepackage{hyperref}
\usepackage{parskip}

\setlength{\parindent}{0pt}
\setlength{\parskip}{8pt}

\pagestyle{empty}

\begin{document}

%====================
% Header
%====================
\begin{center}
    {\LARGE \textbf{Gautam Velpula}} \\[6pt]
    \href{mailto:gautamvelpula@gmail.com}{gautamvelpula@gmail.com} \;|\;
    \href{tel:6785390154}{678-539-0154} \;|\;
    \href{https://www.linkedin.com/in/gautamvelpula}{linkedin.com/in/gautamvelpula}
\end{center}

\vspace{14pt}

\today

\vspace{8pt}

Dear Hiring Team,

I'm excited about the Head of Engineering role at Copilot. I'm an engineering leader who has never stopped being an engineer. I'm at my best when I'm close to the code, close to the engineers building the product, and close to the business problems we're solving together. Copilot's stage and ambition feel like the right place to do exactly that.

At GoTo Foods, my title is Sr Director, but my day-to-day has always been shaped by what the team and the product need. When our subscription service hit a critical failure during a legacy system migration, I didn't delegate the fix---I built the stop-gap service myself. When we needed to consolidate 7 brand apps onto a shared platform, I didn't just draw architecture diagrams---I was in the codebase with our engineers, working through the service contracts, reviewing PRs, and making sure the patterns we chose would actually hold up. Our stack is Node.js, TypeScript, NestJS, Next.js, React Native, and Firebase---the same technologies Copilot is built on---so this isn't theoretical fluency. It's daily practice.

That hands-on credibility is what makes me effective on the other side of the table too. When I sit with product managers, designers, or business stakeholders, I can give honest answers about what's feasible, what's risky, and what we can ship this cycle versus next. I don't need to go back and check with engineering---I \textit{am} engineering. That directness builds trust, speeds up decisions, and keeps roadmap conversations grounded in reality. At Slalom, this was my entire role: translating business strategy into engineering execution for Fortune 500 clients, earning repeat business by being the person both sides trusted to get it right.

What draws me to Copilot specifically is the opportunity to contribute to real growth. I've spent the last chapter of my career inside a large enterprise, and while the work has been meaningful, I'm looking for a place where engineering leadership has a direct, visible impact on the trajectory of the company. Copilot is at the stage where the right engineering decisions---how you scale your platform, how you adopt AI-native workflows, how you grow and retain your team---compound into outsized outcomes. That's the work I want to do.

On AI specifically: I'm not just an advocate, I'm a practitioner. I've driven adoption of Claude Code and Cursor across our engineering organization, defined the architectural guardrails that make LLM-assisted development safe and productive, and built agentic AI capabilities into our consumer products. I understand both the leverage and the risks, and I know how to help a team move fast with these tools without sacrificing the quality and reliability that a fintech product demands.

I'd welcome the chance to talk about how I can help Copilot's engineering team scale---as a leader who builds culture and process, and as an engineer who still ships.

\vspace{16pt}

Sincerely, \\
Gautam Velpula

\end{document}